% Generated draft: requires mathematical and source review.
\documentclass[11pt]{article}
\usepackage[margin=25mm]{geometry}
\usepackage{fontspec}
\setmainfont{Latin Modern Roman}
\usepackage{amsmath,amssymb,mathtools,unicode-math}
\setmathfont{Latin Modern Math}
\usepackage{enumitem,needspace,xurl,hyperref}
\setlength{\parindent}{0pt}
\setlength{\parskip}{5pt}
\setlength{\emergencystretch}{3em}
\newcommand{\N}{\mathbb{N}}
\newcommand{\Z}{\mathbb{Z}}
\newcommand{\Q}{\mathbb{Q}}
\newcommand{\R}{\mathbb{R}}
\newcommand{\C}{\mathbb{C}}
\newcommand{\F}{\mathbb{F}}
\newcommand{\A}{\mathbb{A}}
\newcommand{\PP}{\mathbb{P}}
\newcommand{\E}{\mathbb{E}}
\newcommand{\eps}{\varepsilon}
\newcommand{\dd}{\,\mathrm{d}}
\begin{document}
\section*{UnsolvedMath ID: 30000803}
\textbf{Source problem code:} catalogue entry
\subsection{Definitions and mathematical statement}
Beyond curvelets and shearlets, are there other constructions based on operations on the plane that guarantee reproduction of a signal while offering more flexible analytic and geometric properties?

The source asks for alternative signal-representation constructions with an exact reproduction guarantee and additional analytic/geometric flexibility. It does not specify the signal space, generating functions, allowed planar operations, parameter set, coefficient rule, reconstruction topology, or the properties to be varied. Translations, rotations, anisotropic dilations, and shears are examples of relevant geometric operations, not a source-prescribed list.

One possible mathematical framework would choose a Hilbert space of planar signals and analysis/synthesis families with
\[
 f=\sum_j\langle f,\psi_j\rangle\widetilde\psi_j
\]
for every signal in that space, in a specified convergence sense. A continuous transform would instead require an appropriate integral reproduction formula. The record does not select a discrete frame, basis, continuous transform, or another mechanism, so this display only explains the missing reproduction criterion; it is not imposed as the definition of an acceptable answer.

The phrase ``more flexible analytic and geometric properties'' gives no explicit constraints, benchmark, optimality criterion, or tradeoffs relative to curvelets and shearlets. The record also does not define when a construction counts as genuinely different from those families rather than a variant. No chosen support, smoothness, sparsity, directional, or adaptation theorem is inserted.

The original is an existence/construction research direction. The assessment mentions the lack of a complete classification, but the source question does not explicitly demand a classification of every possible construction. That broader task is not silently substituted for the requested alternatives.
\subsection{Short English statement}
The question seeks new planar geometric signal systems with guaranteed reconstruction and more flexible properties than the cited families. The reproduction framework and desired flexibility are unspecified, so it remains a construction direction rather than a precise theorem or classification demand.
\subsection{Sources}
\begin{itemize}
\item Local UnsolvedMath registry, permanent ID 30000803, code catalogue entry, ``Flexible Geometric Signal-Reproduction Systems.'' Proposed year 2007; proposer unspecified.
\item \emph{Wavelet and Multiscale Methods} (2008), Oberwolfach Reports, pp.\ 2079--2172, DOI 10.4171/owr/2007/36; \url{https://doi.org/10.4171/owr/2007/36}. The displayed report year 2008 differs from the DOI and proposed year 2007.
\item Supplied ``adaptive MRA and shearlet constructions'': \url{https://doi.org/10.1137/090780912}.
\item Statement status exact; the supplied verification calls the construction question clear and complete and says it was retained byte-for-byte. Its desired output properties still need specification.
\item Literature triage checked 2026-08-21: open. It describes a broad request for reproducible directional constructions and says existing curvelet, shearlet, and adaptive-MRA work does not give a complete classification. This assessment does not establish a precise nonexistence or classification target. Created 2026-08-25; updated 2026-09-16. Current status and source scope require independent verification.
\item Registry curation by Ulam AI / UnsolvedMath Contributors, CC BY 4.0; underlying report and literature provenance remain as supplied. Unreviewed statement-only draft.
\end{itemize}
\end{document}
